\section*{Revision Direction: External Benchmarks, External Validation, and Supporting-Artifact Upgrade Plan}

\subsection*{Overall Revision Principle}
The next revision stage should prioritize \emph{protocol-compatible external validation} rather than adding a large number of disconnected benchmarks. The core principle is to preserve the current CPPB evaluation logic, matched-setting comparison protocol, and evidence-bounded claim scope, while gradually extending the manuscript from controlled evidence to externally checkable evidence. In practical terms, the next round of revision should focus on three coordinated goals:
\begin{enumerate}
    \item adding one or two protocol-compatible external benchmarks that can be mapped into the current PER/AC/TSR-style evaluation pipeline;
    \item strengthening the external practical baseline family under the same downstream setting;
    \item upgrading the remaining high-value supporting tables from manuscript-level controlled slices into CSV-backed, script-filled, and reproducible appendix artifacts.
\end{enumerate}

\subsection*{Priority 1: Add a Protocol-Compatible Text Anonymization Benchmark}
The highest-priority external benchmark candidate is a \emph{text anonymization benchmark} that is naturally compatible with the current span-level privacy mediation setup. The best immediate target is a benchmark of the TAB family. This direction is attractive because it is already aligned with anonymization-oriented evaluation rather than generic NER alone. Therefore, it can be adapted into the present paper with relatively limited protocol drift.

The recommended revision strategy is:
\begin{enumerate}
    \item convert benchmark instances into prompt-style inputs;
    \item preserve the current matched comparison setting across raw, regex-only, NER-only, generic de-identification, enterprise staged redaction, and the proposed method;
    \item report direct-exposure and utility-preservation results under the same output protocol as CPPB;
    \item use this benchmark primarily as the first external transfer slice for text-only validation.
\end{enumerate}

This benchmark should be used first to support three manuscript goals:
\begin{enumerate}
    \item stronger external transfer beyond CPPB;
    \item stronger external practical baseline comparison;
    \item a more credible category-wise supporting analysis.
\end{enumerate}

\subsection*{Priority 2: Add a Clinical PHI Benchmark for External Generalization}
The next strongest external validation direction is a \emph{clinical de-identification benchmark}, especially one in the i2b2-style PHI benchmark family. This is not as directly agent-oriented as CPPB, but it provides a strong external domain where privacy-sensitive spans are real, context-dependent, and widely recognized by the community.

The manuscript should not claim that such a benchmark is inherently agent-compatible. Instead, the revision should explicitly state that:
\begin{enumerate}
    \item the benchmark is used as an external de-identification transfer setting;
    \item benchmark samples are wrapped into prompt-style tasks so that the same mediation pipeline can be evaluated under matched conditions;
    \item the goal is external validation of span protection and utility preservation, not a claim of native agent-workflow equivalence.
\end{enumerate}

This benchmark is especially valuable for strengthening the medical-content story already present in the paper and for reducing the risk that the current evaluation is perceived as CPPB-only.

\subsection*{Priority 3: Add OCR/Document Benchmarks for the Multimodal Slice}
To improve the current multimodal appendix slice, the revision should move from manuscript-level controlled multimodal evidence toward at least one \emph{public OCR/document benchmark}. The most useful candidates are datasets in the form-understanding, receipt-understanding, or document-VQA family.

A practical strategy is to separate the multimodal gap into two sub-problems:
\begin{enumerate}
    \item \textbf{OCR/document extraction robustness:} use a noisy scanned-document or receipt-style benchmark to test OCR-mediated sensitive span extraction and sanitization under layout noise, OCR noise, and semi-structured documents;
    \item \textbf{Downstream utility preservation:} use a document question-answering benchmark to test whether the protected prompt/document representation still supports useful downstream answers.
\end{enumerate}

Under this design, the multimodal revision does not need to be a perfect one-shot benchmark package. It is sufficient to:
\begin{enumerate}
    \item choose one public document/OCR benchmark for extraction-focused validation;
    \item optionally choose one public document-QA benchmark for utility-focused validation;
    \item convert both into the existing protected-vs-raw matched evaluation format;
    \item report them honestly as external multimodal supporting validation rather than as a replacement for CPPB.
\end{enumerate}

\subsection*{Priority 4: Use an Agent Security Environment for Boundary-Misuse Validation}
The current paper’s distinctive contribution is not only span masking, but also \emph{propagation control across retrieval, memory, and tool boundaries}. Therefore, a separate external validation track should be considered for \emph{boundary misuse, restoration misuse, and tool-stage leakage}. For this purpose, a benchmark or environment in the agent security / tool-using agent family is useful.

This direction should not be presented as a replacement for anonymization benchmarks. Instead, it should be framed as:
\begin{enumerate}
    \item an external robustness environment for the threat model;
    \item a way to test restoration-trigger attacks, unsafe early restoration, and boundary-policy bypass attempts;
    \item a complementary validation track for BLR/RSR-style claims.
\end{enumerate}

In short, anonymization-style external benchmarks should validate \emph{text privacy mediation}, while agent-security-style environments should validate \emph{boundary control and restoration policy robustness}.

\subsection*{Revision Direction for the External Practical Baseline Family}
The current manuscript still needs a stronger external practical baseline family. The recommended strategy is incremental rather than open-ended. The next revision should not add many unrelated baselines at once. Instead, it should proceed in the following order:
\begin{enumerate}
    \item keep the current enterprise staged redaction comparator as the internal practical baseline;
    \item add one engineering-recognizable external de-identification pipeline under matched CPPB-style evaluation;
    \item add one prompted LLM zero-shot de-identification baseline with a fixed prompt template and frozen downstream setting;
    \item only after these are stable, consider additional enterprise-grade or more generic external pipelines.
\end{enumerate}

The key requirement is that all new baselines must be evaluated under:
\begin{enumerate}
    \item the same prompt slice;
    \item the same downstream backend condition;
    \item the same metric family;
    \item the same evidence boundary.
\end{enumerate}

This prevents benchmark drift and avoids introducing comparisons that look stronger superficially but are not truly matched.

\subsection*{Revision Direction for the Remaining Controlled Supporting Tables}
The manuscript still contains several supporting tables that are not yet fully regenerated. These should not all be upgraded at once. A staged priority order is recommended.

\subsubsection*{Stage A: High-Value Supporting Slices}
The first upgrade targets should be the appendix slices that most directly affect reviewer confidence in the central claim:
\begin{enumerate}
    \item multimodal evaluation;
    \item hard-case robustness analysis;
    \item category-wise analysis.
\end{enumerate}

These slices are more valuable than lower-priority appendix material because they explain \emph{where the method succeeds, where it degrades, and what remains risky}. Each should be upgraded into:
\begin{enumerate}
    \item a CSV-backed record file;
    \item a deterministic fill-table script;
    \item a clear status note in the appendix reproducibility map.
\end{enumerate}

\subsubsection*{Stage B: Cross-Model Supporting Slice}
The cross-model table is important, but its main gap is not the absence of a new benchmark. Its main gap is the lack of exact provenance metadata. Therefore, the next revision should prioritize:
\begin{enumerate}
    \item exact backbone identifiers;
    \item decoding settings;
    \item configuration logs;
    \item rerun manifests.
\end{enumerate}

This means the cross-model slice should be upgraded primarily through \emph{metadata completion and rerun traceability}, not through benchmark expansion alone.

\subsection*{Revision Direction for External Transfer and External Validation}
The manuscript should explicitly present external transfer as a \emph{staged validation roadmap}. A recommended three-step structure is:
\begin{enumerate}
    \item \textbf{Text-only transfer:} add one anonymization-oriented external benchmark;
    \item \textbf{Domain transfer:} add one clinical PHI benchmark;
    \item \textbf{Multimodal transfer:} add one OCR/document benchmark and optionally one document-QA benchmark.
\end{enumerate}

If the revision cycle can only support two external additions, the recommended pair is:
\begin{enumerate}
    \item one anonymization-oriented text benchmark;
    \item one OCR/document or document-QA benchmark.
\end{enumerate}

This gives the paper the strongest improvement per unit effort while remaining aligned with the current story.

\subsection*{Revision Direction for Reproducibility Metadata}
Several remaining weaknesses are primarily \emph{metadata completeness problems}, not scientific-claim problems. These should be addressed through explicit manifest generation rather than prose-only acknowledgement. The revision should create separate reproducibility manifests for:
\begin{enumerate}
    \item cross-model backbone identifiers and decoding settings;
    \item multimodal OCR engine, version, and preprocessing pipeline;
    \item latency host hardware, concurrency condition, and prompt-length distribution.
\end{enumerate}

The manuscript should then reference these manifests explicitly in the appendix reproducibility map. This is preferable to leaving such information only in descriptive text, because it turns a prose limitation into a checkable artifact boundary.

\subsection*{Revision Direction for the CPPB Data Card}
The CPPB data card can be strengthened further without changing the core benchmark. The next revision should expand the data-card section in the following order:
\begin{enumerate}
    \item source-level release metadata;
    \item clearer annotation examples;
    \item fuller split and label semantics;
    \item multimodal provenance and licensing notes;
    \item release-scope statement aligned with the appendix reproducibility section.
\end{enumerate}

The goal is not to oversell CPPB as a mature community benchmark. The goal is to make the benchmark documentation sufficiently explicit that the paper’s evidence boundary is transparent and reviewer-facing concerns about hidden assumptions are reduced.

\subsection*{Recommended Execution Order}
A practical execution order for the next revision cycle is:
\begin{enumerate}
    \item add one external text anonymization benchmark;
    \item add one stronger external practical baseline under the same protocol;
    \item upgrade the multimodal, hard-case, and category-wise tables to repo-backed supporting artifacts;
    \item complete cross-model and multimodal metadata manifests;
    \item add one external OCR/document validation slice;
    \item optionally add one agent-security robustness environment for boundary-misuse validation.
\end{enumerate}

\subsection*{Suggested Paper-Level Positioning After Revision}
After the above modifications, the manuscript can be positioned more strongly as follows:
\begin{quote}
The paper provides controlled systems evidence on CPPB, extends the validation with protocol-compatible external anonymization and document benchmarks, strengthens the external practical baseline family under matched conditions, and upgrades the main supporting appendix slices into reproducible artifacts while preserving an evidence-bounded claim scope.
\end{quote}

\subsection*{Concise Summary for the Next Revision Round}
In summary, the next revision should not aim for uncontrolled benchmark expansion. Instead, it should follow a disciplined external-validation pathway:
\begin{enumerate}
    \item first strengthen text-side external transfer;
    \item then strengthen multimodal/OCR validation;
    \item then strengthen boundary-robustness validation;
    \item and in parallel convert the highest-value remaining controlled slices into reproducible supporting artifacts with explicit manifests.
\end{enumerate}

This strategy is the most efficient way to improve reviewer confidence without distorting the current architecture, evaluation protocol, or evidence scope of the manuscript.